% Threats to Validity — 4-quadrant table (Wohlin et al. 2012 taxonomy)
% Hand-curated 2026-05-17, designed to be paste-ready into the paper.
\section{Threats to Validity}
\label{sec:threats}

We organize threats following the classical taxonomy of Wohlin
et~al.~\cite{wohlin2012experimentation} into four categories: construct,
internal, external, and conclusion validity. For each threat we describe the
risk and the mitigation already applied or its residual scope.

\begin{table}[t]
\centering
\caption{Threats to validity and applied mitigations.}
\label{tab:threats}
\small
\begin{tabular}{p{0.10\linewidth} p{0.36\linewidth} p{0.46\linewidth}}
\hline
\textbf{Category} & \textbf{Threat} & \textbf{Mitigation / Residual scope} \\
\hline
\multirow{3}{*}{Construct}
 & T1: \emph{Flakiness} is operationalized as ``the same test exhibits both
   pass and fail outcomes across repeated runs of the same SUT version.''
   This excludes intermittent failures masked by retry policies.
 & We disable all in-test retry and report the raw first-outcome per run
   (\S\ref{sec:harness}). Suite-level outcomes are recomputed offline from
   per-test CSVs to avoid harness-side reattribution. \\
 & T2: ``Database isolation'' is reified only at the PostgreSQL layer.
   Other shared state (object storage, message queues, in-cluster caches)
   is preview-scoped but not formally bounded.
 & The five subjects exercise only the PostgreSQL state path under test.
   We document the architectural boundary in \S\ref{sec:operator} and flag
   it as a scope decision rather than a measurement defect. \\
 & T3: ``Cycle time'' aggregates operator-side step durations
   (smoke + regression + e2e + restore overhead) but excludes
   client-side waits.
 & Wall-clock of the controlling \texttt{Preview} object envelope is logged
   independently for cross-check; the two series differ by at most 3\% across
   1{,}200+ runs (see Table~\ref{tab:rq3}). \\
 & T4: For subject S5 (PetClinic), the RQ3 \texttt{mode=restore} cycle-time
   measurements were collected with SUT image \texttt{v3.4.0} (2026-05-16),
   whereas RQ1, RQ2, RQ3-baseline, and RQ5 measurements use the corrected
   image \texttt{v3.4.0-fix7} (2026-05-17).
 & RQ3 cycle-time is measured at operator-side phase boundaries
   (\texttt{postgres-migrate}, \texttt{checkpoint-restore},
   \texttt{checkpoint\_total}) which are image-independent: the Spring~Boot
   SUT container starts identically across the two image versions, and the
   tests do not affect the measured restore/migration phase durations. The
   S5 RQ3 envelope (14.2\,s checkpoint, 87\,s pipeline on) is therefore
   directly comparable to the other 4 subjects. The image-version drift is
   documented in \texttt{HARNESS\_FIXES.md} for reproducibility. \\
\hline
\multirow{3}{*}{Internal}
 & I1: Selection of subjects (S1--S5) reflects the openly available
   PostgreSQL-backed Kubernetes-deployable applications at the time of study.
   This is a convenience sample.
 & We deliberately span four runtimes (Python/Flask, Go, Python/Django, TS/
   Prisma, Java/Spring) and three persistence patterns (raw SQL, GORM,
   Django ORM, Prisma, Spring~Data) to reduce single-stack bias. \\
 & I2: The baseline comparator (\texttt{mode=migration}) shares the same
   operator with the treatment (\texttt{mode=restore}). Effects ascribed
   to the isolation mechanism could conflate with operator-side artifacts.
 & The two modes share \emph{all} operator code paths
   except the suite-restore step, where one branch invokes \texttt{pg\_restore}
   and the other re-runs the migration job (commit \texttt{b956e4a}, file
   \texttt{checkpoint.go}). Both modes share the same reconciler,
   eventing, and failure-handling code paths; the residual operator-side
   confound is therefore bounded to ${\sim}30$ lines of dispatch logic. \\
 & I3: The harness re-uses run-id seeds across subjects; a defect in
   seed handling could correlate runs spuriously.
 & Seeds are derived from \texttt{(subject, condition, iteration)} tuples
   and verified independent via \texttt{analysis/check\_k\_consistency.py},
   which confirms 100\% completion across all $K \in \{1,2,4,8\}$ batches. \\
\hline
\multirow{3}{*}{External}
 & E1: All experiments execute on a single Azure Kubernetes Service cluster
   (3$\times$\texttt{D4s\_v3}, 12 vCPU / 48~GiB).
 & The findings are framed as \emph{lower bounds} on the speed-up and
   \emph{upper bounds} on residual flakiness; we explicitly do not claim
   generalization to GKE, EKS, on-prem K8s, or single-node Kind. A Kind
   replication on S1--S3 ($N=20$ per condition) is listed as
   immediate follow-up work in \S\ref{sec:future}. \\
 & E2: The operator under study is a $\sim$12k LoC custom implementation;
   findings may not transfer to vCluster, Argo Rollouts, Tilt, or
   commercial ephemeral-env tools.
 & The mechanism (\emph{pre-warmed checkpoint $\rightarrow$ pg\_restore})
   is decoupled from the operator's reconciler design and could be
   implemented by any controller exposing the same \texttt{Preview} CRD shape.
   The comparison table in
   Table~\ref{tab:related-work-comparison} positions our claims against the
   nearest neighbours in the design space. \\
 & E3: Subject SUTs are open-source benchmarks; production workloads
   may exhibit different failure profiles.
 & We classify failure causes into 8 assertion categories
   (\S\ref{sec:assertion-categories}); the isolation-probe family
   (\texttt{run\_log\_clean}, \texttt{entity\_count\_matches\_seed})
   constitutes 38\% of total assertions and is independent of
   application-level semantics. \\
\hline
\multirow{3}{*}{Conclusion}
 & C1: Statistical power could be insufficient to detect small effects.
 & Post-hoc power analysis (Table~\ref{tab:T1_effect_sizes}) reports
   observed power $\ge 0.99$ for all RQ1 and RQ3 contrasts at
   $\alpha=0.05$, with Cohen's $d \in [2.4, 20.9]$ for cycle time and
   $|h|=\pi$ (theoretical maximum) for flakiness elimination. Cliff's
   $\delta = \pm 1.0$ confirms complete separation in the
   non-parametric setting. \\
 & C2: Multiple comparisons across five subjects and four RQs inflate
   the family-wise error rate.
 & Holm--Bonferroni correction is applied within each RQ family
   (\S\ref{sec:analysis-methodology}). Effect sizes report bootstrap 95\%
   CIs ($B=10{,}000$, seed=20260517; Table~\ref{tab:T1_bootstrap_ci}) to
   shift inference from $p$-values onto magnitudes. \\
 & C3: The ``zero-variance'' claim for RQ1 is degenerate
   ($\mathrm{Var}=0$ in both conditions for some subjects), inviting
   misreading.
 & We report variance equality via Bartlett's test
   (Table~\ref{tab:T1_variance_tests}) and clarify that
   $\mathrm{Var}(\text{iso}=\text{True})=0$ is the load-bearing finding;
   $\mathrm{Var}(\text{iso}=\text{False})$ being also $\approx 0$ for some
   subjects reflects \emph{deterministic} failure modes triggered by
   shared state, not the absence of flakiness (the failing
   \emph{test names} vary; the binary indicator does not). \\
\hline
\end{tabular}
\end{table}
